
\documentclass{article}
\usepackage{colortbl}
\usepackage{makecell}
\usepackage{multirow}
\usepackage{supertabular}

\begin{document}

\newcounter{utterance}

\centering \large Interaction Transcript for game `eqbench', experiment `eq\_bench', episode 83 with reflexion{-}Qwen3.5{-}27B.
\vspace{24pt}

{ \footnotesize  \setcounter{utterance}{1}
\setlength{\tabcolsep}{0pt}
\begin{supertabular}{c@{$\;$}|p{.15\linewidth}@{}p{.15\linewidth}p{.15\linewidth}p{.15\linewidth}p{.15\linewidth}p{.15\linewidth}}
    \# & \multicolumn{2}{c}{Player} && \multicolumn{2}{c}{Game Master} \\
    \hline

    \theutterance \stepcounter{utterance}  
    & & & \multicolumn{4}{p{0.6\linewidth}}{
        \cellcolor[rgb]{0.9,0.9,0.9}{
            \makecell[{{p{\linewidth}}}]{
                \texttt{\tiny{[P1$\langle$GM]}}
                \texttt{You have attempted this same task before and did not succeed. These are the plans you wrote after those attempts.} \\
\texttt{Trial \#1: New plan: My previous attempt failed because I assigned a high score (7) to "Triumphant," which contradicts the context of the dialogue. Mari's final line ("Better than drinking myself into oblivion") is a sharp, defensive retort to a personal attack, not a declaration of victory. While she is clearly "Unapologetic" (high score) and "Defensive" (moderate score), she is not "Triumphant" as the conflict remains unresolved and hostile. I also underestimated "Defensive" given the aggressive nature of Charlie's accusations. Next time, I will strictly align the scores with the specific emotional tone of the final exchange: keep "Unapologetic" high, increase "Defensive" to reflect the reaction to the insult, and set "Triumphant" to 0 or very low, as there is no evidence of her feeling victorious. I will ensure the scores reflect a realistic emotional profile where the character is standing her ground but not celebrating a win.} \\
\\ 
\texttt{Trial \#2: New plan: My previous attempt failed because I assigned a non{-}zero score (1) to "Triumphant," which misinterprets the hostile and unresolved nature of the final exchange. Mari's retort ("Better than drinking myself into oblivion") is a sharp, defensive counter{-}attack to a personal insult, not a declaration of victory. While she is clearly "Unapologetic" (high score) and "Defensive" (high score), she is not "Triumphant" as the conflict remains tense and she is merely defending her integrity against Charlie's accusations. I also need to ensure "Guilty" remains 0, as she shows no remorse. Next time, I will strictly align the scores with the specific emotional tone: keep "Unapologetic" and "Defensive" high (e.g., 9 and 8), and explicitly set "Triumphant" to 0 to reflect the lack of victory, ensuring the emotional profile accurately depicts a character standing her ground in a heated argument rather than celebrating a win.} \\
\\ 
\\ 
\\ 
\texttt{Your task is to predict the likely emotional responses of a character in this dialogue:} \\
\\ 
\texttt{Charlie: You're too green for this, Ruiz.} \\
\texttt{Mari: And you're too washed up, Redburn.} \\
\texttt{Charlie: I've been on this beat longer than you've been alive, kid.} \\
\texttt{Mari: And yet, here we are, stuck on the same case. Seems like experience isn't everything.} \\
\texttt{Charlie: Don't get smart with me, Mari.} \\
\texttt{Mari: Someone has to be the smart one here.} \\
\texttt{Charlie: You think I don't know what you're doing? Trying to prove yourself, make a name?} \\
\texttt{Mari: Better than drinking myself into oblivion.} \\
\texttt{Charlie: That's a low blow, Ruiz.} \\
\texttt{{[}End dialogue{]}} \\
\\ 
\texttt{At the end of this dialogue, Mari would feel...} \\
\texttt{Defensive} \\
\texttt{Unapologetic} \\
\texttt{Guilty} \\
\texttt{Triumphant} \\
\\ 
\texttt{Give each of these possible emotions a score from 0{-}10 for the relative intensity that they are likely to be feeling each.} \\
\\ 
\texttt{You must output in the following format, including headings (of course, you should give your own scores), with no additional commentary:} \\
\\ 
\texttt{Defensive: <score>} \\
\texttt{Unapologetic: <score>} \\
\texttt{Guilty: <score>} \\
\texttt{Triumphant: <score>} \\
\\ 
\\ 
\texttt{{[}End of answer{]}} \\
\\ 
\texttt{Remember: zero is a valid score, meaning they are likely not feeling that emotion. You must score at least one emotion > 0.} \\
\\ 
\texttt{Your answer:} \\
            }
        }
    }
     \\ \\

    \theutterance \stepcounter{utterance}  
    & \multicolumn{4}{p{0.6\linewidth}}{
        \cellcolor[rgb]{0.561,0.737,0.561}{
            \makecell[{{p{\linewidth}}}]{
                \texttt{\tiny{[P1$\rangle$GM]}}
                \texttt{Defensive: 9} \\
\texttt{Unapologetic: 9} \\
\texttt{Guilty: 0} \\
\texttt{Triumphant: 0} \\
            }
        }
    }
    & & \\ \\

    \theutterance \stepcounter{utterance}  
    & & & \multicolumn{2}{p{0.3\linewidth}}{
        \cellcolor[rgb]{0.9,0.9,0.9}{
            \makecell[{{p{\linewidth}}}]{
                \texttt{\tiny{[GM$|$GM]}}
                \texttt{Defensive: 9} \\
\texttt{Unapologetic: 9} \\
\texttt{Guilty: 0} \\
\texttt{Triumphant: 0} \\
            }
        }
    }
    & & \\ \\

    \theutterance \stepcounter{utterance}  
    & & & \multicolumn{2}{p{0.3\linewidth}}{
        \cellcolor[rgb]{0.9,0.9,0.9}{
            \makecell[{{p{\linewidth}}}]{
                \texttt{\tiny{[GM$|$GM]}}
                \texttt{Defensive: 4} \\
\texttt{Unapologetic: 8} \\
\texttt{Guilty: 0} \\
\texttt{Triumphant: 6} \\
            }
        }
    }
    & & \\ \\

    \theutterance \stepcounter{utterance}  
    & & & \multicolumn{2}{p{0.3\linewidth}}{
        \cellcolor[rgb]{0.9,0.9,0.9}{
            \makecell[{{p{\linewidth}}}]{
                \texttt{\tiny{[GM$|$GM]}}
                \texttt{game\_result = LOSE} \\
            }
        }
    }
    & & \\ \\

    \theutterance \stepcounter{utterance}  
    & & \multicolumn{4}{p{0.6\linewidth}}{
        \cellcolor[rgb]{0.95,0.95,0.95}{
            \makecell[{{p{\linewidth}}}]{
                \texttt{\tiny{[P1$\leftrightarrow$P]}}
                \texttt{Post{-}episode Reflexion (not sent to GM):} \\
\texttt{New plan: My previous attempt failed because I assigned identical high scores (9) to both "Defensive" and "Unapologetic," failing to distinguish the primary emotional driver. While Mari is certainly unapologetic, her final line ("Better than drinking myself into oblivion") is a direct, stinging reaction to Charlie's personal attack ("That's a low blow"), making "Defensive" the dominant, most intense emotion. "Unapologetic" is a secondary trait of her stance, not the immediate reaction to the insult. Next time, I will differentiate the scores to reflect this nuance: assign the highest score (e.g., 10) to "Defensive" to capture the intensity of the reaction to the low blow, assign a slightly lower but still high score (e.g., 8) to "Unapologetic" to show her lack of regret, and keep "Guilty" and "Triumphant" at 0. This ensures the scoring hierarchy accurately reflects the dialogue's climax.} \\
            }
        }
    }
    & \\ \\

\end{supertabular}
}

\end{document}
